\documentclass[10pt,a4paper]{amsart}
\usepackage[margin=19mm]{geometry}
\usepackage{amsmath,amssymb,mathtools}
\usepackage[scr=rsfs,cal=euler]{mathalfa}
\usepackage{xfrac}
\usepackage[unicode,hidelinks,bookmarks=false]{hyperref}
\newcommand{\ie}{i.e.\ }
\newcommand{\cf}{cf.\ }
\newcommand{\resp}{respectively\ }
\newcommand{\Subsection}{\subsection}
\providecommand{\nobreakdash}{\nobreak\hbox{-}}
\setlength{\emergencystretch}{2em}
\multlinegap0pt
\usepackage{bm}
\usepackage{bbm}
\usepackage{dsfont}
\usepackage{xspace}
\usepackage{cancel}
\usepackage{mathtools}
\usepackage{amsthm}
\makeatletter
\@ifundefined{theorem}{\newtheorem{theorem}{Theorem}}{}
\@ifundefined{lemma}{\newtheorem{lemma}[theorem]{Lemma}}{}
\makeatother
\newcommand{\PFk}{P_{F_k}}
\newcommand{\R}{\mathbb R}
\title[Harmonic Maps into Euclidean Buildings and Non-Archimedean S]{Harmonic Maps into Euclidean Buildings and Non-Archimedean Superrigidity}
\author{Christine Breiner, Ben K. Dees, Chikako Mese}
\date{Source: arXiv:2408.02783v2 (CC BY 4.0)}
\begin{document}
\maketitle

\noindent\emph{Source and licence.} Harmonic Maps into Euclidean Buildings and Non-Archimedean Superrigidity. Version arXiv:2408.02783v2 (CC BY 4.0), distributed under the Creative Commons Attribution 4.0 International licence, as recorded in the arXiv licence metadata for this exact version and shown on the abstract page https://arxiv.org/abs/2408.02783v2. Full rights notice: licenses/arxiv-2408.02783-CC-BY-4.0.txt.\par\smallskip

\noindent\textit{Editorial scope.} Counted unit: the Lemma whose source label is \texttt{omega*} in the article (source0.tex lines 1468--1477, source label \texttt{omega*}), delivered together with the article's own complete original proof of it (source0.tex lines 1479--1560, one \texttt{proof} environment of 82 lines). No mathematics is written, repaired or re-derived: statement and proof are the article's own text line for line. Required context retained uncounted: upsilonprime (lines 1395--1404). Every definition, display and label of those lines is retained unchanged. Omitted from this package and not redistributed: 1--1394, 1405--1467, 1478--1478, 1561--2324. Nothing inside a retained excerpt is omitted or altered: every retained body is delivered exactly as the article's lines are written, so no omission receipt of any kind accompanies this package. Citations: the retained excerpts carry no citation command at all, so no bibliography entry of the article is transcribed and no reference is redistributed. Reference adaptation: none was needed and none was made; every reference inside the delivered unit resolves inside it, so no unresolved reference or citation is produced. Printed slips kept literal: no printed word, symbol, hypothesis or number of the counted statement or of its proof was altered. Layout and class: the article's own class is \texttt{amsart} with packages including amsmath, tikz-cd, mathtools, amsfonts, graphicx, verbatim, textcomp, amssymb, cite, color, xy, hyperref, mathrsfs, soul; the wrapper is the lane's pinned \texttt{amsart} editorial wrapper at 10pt on A4, which restates the theorem-like environments the retained text needs and adds the article's own macro definitions that the retained text needs (\texttt{\textbackslash PFk}, \texttt{\textbackslash R}), with extra wrapper packages (bm, bbm, dsfont, xspace, cancel, mathtools). The delivered numbering is the unit's own. Assets: no retained span contains a figure environment, a graphics-inclusion command, a PDF-page inclusion, a table or a TikZ picture; no image file is copied into this package, the package declares no asset dependency and no image file is redistributed with it. Licence: version arXiv:2408.02783v2 of this article is distributed under the Creative Commons Attribution 4.0 International licence, as recorded in the arXiv licence metadata for this exact version and shown on the abstract page https://arxiv.org/abs/2408.02783v2. A full rights notice is delivered at licenses/arxiv-2408.02783-CC-BY-4.0.txt. Characters: every retained excerpt is pure ASCII, so the delivered engine, XeLaTeX with its default Latin Modern text font, reports no missing character.
\begin{lemma} \label{upsilonprime}
Fix $\tau >0$ and $i=1, \dots, m$.  
\[
\mbox{If }
\Upsilon_k(\tau, i):=\left\{ p \in  B_{r_0}(0):   (1+ \tau)^2 \lambda_i^2  \leq  \left|\frac{\partial u_k}{\partial x^i} \right|^2(p)\right\},
\mbox{ 
then }
\lim_{k \to \infty}\mu_0(\Upsilon_k(\tau, i))=0.
\]
\end{lemma}

\begin{lemma} \label{omega*}
Fix  $\delta>0$ and $i\in \{1, \dots, m\}$. 
\[
\mbox{If }
\Theta_k(\delta, i): =\left\{ p \in B_{r_0}(0):   \left|\frac{\partial u_k^i}{\partial x^i}(p) \right|^2 \leq  (1-\delta)^2\lambda_i^2 \right\}, 
\mbox{
then } 
\lim_{k\to\infty}\mu_0(\Theta_k(\delta, i))=0. 
\]
\end{lemma}

\begin{proof} 
We consider the case $i=1$ as all other cases follow similarly.    On the contrary, assume there exists a subsequence of $k \rightarrow \infty$ (which we will still denote as $k$ by an abuse of notation) such that $\lim_{k \to \infty}\mu_0(\Theta_k(\delta, 1)) \geq \beta >0$.


Let ${\bf B} \subset \R^{n-1}\simeq \{0\} \times \R^{n-1}$ be the ball of radius $r_0$ centered at the origin. For $p \in B_{r_0}(0)$, write $p = (p_1, \bar p)$ where $\bar p:= (p_2, \dots, p_n) \in {\bf B}$. 
Use the projection of $B_{r_0}(0) \rightarrow {\bf B}$, $(p_1,\bar p) \mapsto \bar p$ to view $B_{r_0}(0)$ as a fiber bundle over ${\bf B}$ with intervals as fibers.  
More precisely,  to each $\bar p \in {\bf B}$, we associate an interval  $I_{\bar p}:=(-\rho(\bar p), \rho(\bar p)) \subset \R$  where $\rho(\bar p)=\sqrt{(3/4)^2 - |\bar p|^2}$. 


For $\bar p \in {\bf B}$, define a subset $\theta_k(\bar p)$ of the interval $I_{\bar p}$ by 
\[
\theta_k(\bar p):= \left\{p_1 \in I_{\bar p}:  \left| \frac{\partial u^1_k}{\partial x^1} (p_1,\bar p) \right|\leq (1-\delta)\lambda_1\right\}.
\]
 Define a subset   $\mathcal A_k$ of the base space ${\bf B}$ by 
$$\mathcal A_k:= \{\bar p \in {\bf B}: \mu_0^1(\theta_k(\bar p))> \beta/4\omega\}
$$ (recall  $\mu_0^k$ denotes the $k$-dimensional Lebesgue measure) where $\omega:=\mu_0^{n-1}({\bf B})$. 
Now suppose that $\liminf_{k \to \infty} \mathcal A_k =0$. Then there exists a subsequence (again labeled by $k$) such that 
 $\lim_{k \to \infty}\mu_0^{n-1}(\mathcal A_k^c) =\omega$. (We use superscript $c$ to denote the complement of a set.)
 
 Then Fubini's theorem implies
\begin{align*}
\frac{\beta}{4} &= \lim_{k \to \infty}\mu_0^{n-1}(\mathcal A_k^c) \cdot \frac{\beta}{4\omega} \geq \lim_{k \to \infty} \int_{\mathcal A_k^c} \mu_0^1(\theta_k(\bar p)) d\mu_0^{n-1}\\& =\lim_{k \to \infty} \int_{{\bf B}}\mu_0^1(\theta_k(\bar p)) d\mu_0^{n-1}= \lim_{k \to \infty}\mu_0(\Theta_k(\delta, 1)) \geq \beta,
\end{align*}
a contradiction.  Thus, $\liminf_{k \to \infty}\mu_0^{n-1}(\mathcal A_k) >0$. 

For $\kappa, \tau>0$ to be chosen later, define a subset $\mathcal B_k$ of the base ${\bf B}$ by
\begin{align*}
\mathcal B_k:= \left\{ \bar p \in {\bf B}:\right.& \left.\left| \frac{\partial u^1_k}{\partial x^1} (p_1, \bar p) \right|\leq (1+\tau) \lambda_1 \text{ for } p_1 \in I_{\bar p}\right.\left.  \text{ except on a subset of measure}< \kappa \right\}.
\end{align*} 
Lemma \ref{upsilonprime} implies that $\lim_{k \to \infty}\mu_0^{n-1}(\mathcal B_k) = \omega$. Thus,  there exists $K$  with the property that $\mathcal A_k \cap \mathcal B_k \neq \emptyset$ for all $k> K$.  Choose $\bar p_k \in \mathcal A_k \cap \mathcal B_k$ for each $k$.

Define a subset $b_k$ of the interval $I_{\bar p_{k}}$ by
\[
 b_k:=\left \{t \in I_{\bar p_k}: \left| \frac{\partial u^1_k}{\partial x^1} (t, \bar p_k) \right| \leq (1+\tau) \lambda_1 \right\}.
\]
For ease of notation, let $$
f_k(t):= u_k^1(t, \bar p_k).
$$
 The inequality $|f_k'(t)|^2 \leq \left|\frac{\partial u_k}{\partial  t}\right|^2(t, \bar p_k)$ and the uniform Lipschitz bounds on $u_k$ in $B_{r_0}(0)$ imply there exists $C>0$ such that 
 $$
 |f_k'(t)| \leq C, \ \ \forall t \in I_{\bar p_k}=(-\rho(\bar p_k), \rho(\bar p_k)).$$


Since $\bar p_k \in \mathcal A_k \cap  \mathcal B_k$,
\[
\mu_0^1(\theta_k(\bar p_k)) > \beta/4\omega \ \ \mbox{ and } \ \ 
\mu_0^1(b_k^c)<\kappa.
\]
Therefore, for $k >K$,
\begin{align*}
f_k(\rho(\bar p_k))- f_k(-\rho(\bar p_k))&=
 \int_{-\rho(\bar p_k)} ^{\rho(\bar p_k)}f_k'(t)\, dt
  \nonumber \\
& \leq \int_{-\rho(\bar p_k)} ^{\rho(\bar p_k)}|f_k'(t)|\, dt \nonumber \\
&=\int_{\theta_k(\bar p_k)} |f_k'(t)|\, dt+\int_{\theta_k(\bar p_k)^c \cap b_\kappa} |f_k'(t)|\, dt+\int_{\theta_k(\bar p_k)^c \cap b_\kappa^c} |f_k'(t)|\, dt 
\nonumber \\
& \leq\mu_0^1(\theta_k(\bar p_k))(1-\delta)\lambda_1+ \mu_0^1(\theta_k(\bar p_k)^c) (1+\tau) \lambda_1+ C \mu_0^1(b_\kappa^c)\nonumber \\
& =   2 \rho(\bar p_k)  \lambda_1
-\mu_0^1(\theta_k(\bar p_k))\delta \lambda_1+ \mu_0^1(\theta_k(\bar p_k)^c) \tau \lambda_1+ C \mu_0^1(b_\kappa^c) \nonumber
\\
&< 2 \rho(\bar p_k)  \lambda_1
-\frac{\beta \delta \lambda_1}{4\omega}+ \mu_0^1(\theta_k(\bar p_k)^c) \tau \lambda_1+ C \kappa. \nonumber
\end{align*}
Thus, by choosing $\kappa, \tau>0$ sufficiently small (depending only on $\beta, \delta, \lambda_1, \omega, C$), we conclude
\begin{equation} \label{rhocontra}
f_k(\rho(\bar p_k))- f_k(-\rho(\bar p_k)) <  2 \rho(\bar p_k)  \lambda_1
-\frac{\beta \delta \lambda_1}{8\omega}, \ \ \forall k >K.
\end{equation}

On the other hand, for any $p \in B_{r_0}(0)$, we can view $u_k^1(p)$, $L^1(p)$ as points in $\R \simeq  \R \times \{0,\dots, 0\} \subset \R^N \simeq A_k \subset \PFk$.  Under this identification,
\[
|L^1(p)- u_k^1(p)| = d_k(L^1(p),  u_k^1(p)).
\]
Since $\sup_{B_{r_0}(0)} d_k(L^1(p),u_k^1(p)) 
   \leq 
   \sup_{B_{r_0}(0)} d_k(L_k(p), u_k(p)) < \frac{1}{k}$,
   we conclude
\[
\lim_{k \rightarrow \infty}  \sup_{p \in B_{r_0}(0)}|L^1(p)- u_k^1(p)| = 0.
\]
In particular, this implies $f_k(\rho(\bar p_k))- f_k(-\rho(\bar p_k)) \to 2 \rho(\bar p_k)\lambda_1$ as $k \rightarrow \infty$, contradicting (\ref{rhocontra}).
\end{proof}

\end{document}
